    \documentclass[a4paper,12pt]{article}

    \usepackage[utf8]{inputenc}
    \usepackage[spanish, es-noshorthands]{babel}
    \usepackage{geometry}
    \usepackage{graphicx}
    \usepackage{multirow}
    \usepackage{array}
    \usepackage{fancyhdr}
    \usepackage{lastpage}
    \usepackage{hyperref}
    \usepackage{float}
    \usepackage{caption}
    \usepackage{booktabs}
    \usepackage[table]{xcolor}
    \usepackage[T1]{fontenc}
    \usepackage{tabularx}
    \usepackage{amssymb}
    \usepackage{pifont}
    \usepackage{listings}

    \hypersetup{colorlinks=true,linkcolor=blue,urlcolor=blue}

    \definecolor{SKTPrimary}{HTML}{3F4A4F}
    \definecolor{SKTDark}{HTML}{242B2E}
    \definecolor{SKTAccent}{HTML}{979C9D}
    \definecolor{SKTLight}{HTML}{CDCFD2}
    \definecolor{SKTMid}{HTML}{6E7375}

    \definecolor{rowGET}{HTML}{DFF0D8}
    \definecolor{rowPOST}{HTML}{D9EDF7}
    \definecolor{rowPATCH}{HTML}{FCF8E3}
    \definecolor{rowDELETE}{HTML}{F2DEDE}

    \geometry{
    top=2.5cm,
    bottom=2.5cm,
    left=3cm,
    right=3cm,
    includehead,
    headheight=130pt,
    headsep=1cm
    }

    % ---------------- ENCABEZADO ----------------
    \pagestyle{fancy}
    \fancyhf{}
    \renewcommand{\headrulewidth}{0pt}

    \fancyhead[C]{
    \small
    \setlength{\tabcolsep}{4pt}
    \renewcommand{\arraystretch}{1.1}

    \    \begin{tabular}{|m{3cm}|m{7.5cm}|m{3.8cm}|}
    \hline
    \multirow{5}{*}{
    \parbox[c][3.5cm][c]{2.5cm}{
    \centering
    \includegraphics[width=2.2cm]{SKTLogo.png}
    }
    }
    & \centering\textbf{EMPRESA DE DESARROLLO E INNOVACIÓN}
    & \parbox[c]{\linewidth}{\centering
    \textbf{CÓDIGO:}\\
    MI-DES-MAN-001
    \vspace{3pt}
    } \\ \cline{2-3}

    & \centering\textbf{PROCESO: DESARROLLO DE SOFTWARE}
    & \parbox[c]{\linewidth}{\centering
    \textbf{VERSIÓN:}\\
    1.0
    \vspace{3pt}
    } \\ \cline{2-3}

    & \centering Desarrollo del Sistema LUPSI
    & \parbox[c]{\linewidth}{\centering
    \textbf{ELABORACIÓN:}\\
    07/04/2026
    } \\ \cline{2-3}

    & \centering Especificación de Autenticación y Validación de Identidad
    & \parbox[c]{\linewidth}{\centering
    \textbf{ÚLTIMA REVISIÓN:}\\
    07/04/2026
    } \\ \hline
    \end{tabular}
    }

    \fancyfoot[R]{\small Página \thepage\ de \pageref{LastPage}}

    % --------------------------------------------------
    \begin{document}
    \raggedbottom

    % ==========================================
    \section{Título}
    % ==========================================
    \textbf{Especificación Técnica de Autenticación Centralizada y Validación Algorítmica de Identidad (Módulo 10) para el Ecosistema LUPSI}

    % ==========================================
    \section{Introducción}
    % ==========================================
    El sistema LUPSI prioriza una experiencia de usuario de \textbf{fricción cero} (\textit{Zero Friction}) sin comprometer la integridad de la identidad ni la seguridad de los datos médicos. Este documento detalla la implementación técnica del motor de autenticación basado en Supabase Auth y el servicio de validación de cédula ecuatoriana mediante el algoritmo Módulo 10.

    La autenticación no solo garantiza el acceso, sino que orquesta la sincronización entre la capa de seguridad gestionada de Supabase y el modelo relacional de la base de datos (perfiles y pacientes), asegurando que cada usuario autenticado tenga una identidad válida y verificable desde el primer segundo de su registro.

    % ==========================================
    \section{Objetivo}
    % ==========================================
    Documentar formalmente el funcionamiento interno, los flujos lógicos y las restricciones de seguridad aplicables al proceso de registro, inicio de sesión y validación de identidad en el backend LUPSI, sirviendo como guía de referencia técnica para mantenimiento y auditoría.

    % ==========================================
    \section{Tecnología y Arquitectura}
    % ==========================================

    \begin{itemize}
        \item \textbf{Backend Framework}: NestJS con arquitectura de microservicios lógica.
        \item \textbf{Identity Provider}: Supabase Auth (Admin SDK para operaciones de registro).
        \item \textbf{Validation Logic}: Servicio especializado \texttt{EcuadorianIdValidatorService}.
        \item \textbf{Data Sync}: Sincronización atómica entre \texttt{auth.users} y \texttt{public.profiles / public.patients}.
    \end{itemize}

    % ==========================================
    \section{Validación de Cédula (Algoritmo Módulo 10)}
    % ==========================================

    El servicio \texttt{EcuadorianIdValidatorService} implementa una validación matemática de tres capas para garantizar que el DNI ingresado sea un documento verídico emitido por el Registro Civil del Ecuador.

    \subsection{Reglas de Validación}

    \begin{enumerate}
        \item \textbf{Longitud y Formato}: Debe contener exactamente 10 dígitos numéricos.
        \item \textbf{Código de Provincia}: Los dos primeros dígitos deben estar en el rango [01, 24] o ser el código 30 (Ecuatorianos en el exterior).
        \item \textbf{Tipo de Persona}: El tercer dígito debe ser menor a 6, indicando que el DNI pertenece a una persona natural y no a una entidad jurídica (RUC).
        \item \textbf{Algoritmo de Módulo 10}: Validación del dígito verificador.
    \end{enumerate}

    \subsection{Lógica del Algoritmo}

    El proceso matemático sigue estos pasos:
    \begin{itemize}
        \item Se extraen los primeros 9 dígitos.
        \item Se multiplican por los coeficientes fijos: \texttt{[2, 1, 2, 1, 2, 1, 2, 1, 2]}.
        \item Si el resultado de una multiplicación es $\ge 10$, se le resta 9.
        \item Se suman todos los resultados parciales (\texttt{totalSum}).
        \item Se calcula la decena superior inmediata de \texttt{totalSum}.
        \item El dígito verificador resultante debe ser igual a la diferencia entre la decena superior y la suma total.
    \end{itemize}

    \begin{lstlisting}[language=TypeScript, caption={Extracto de implementación del validador}, basicstyle=\small\ttfamily]
    const coefficients = [2, 1, 2, 1, 2, 1, 2, 1, 2];
    for (let i = 0; i < 9; i++) {
    let product = parseInt(dni.charAt(i)) * coefficients[i];
    if (product >= 10) product -= 9;
    totalSum += product;
    }
    const nextTen = Math.ceil(totalSum / 10) * 10;
    const expectedDigit = (nextTen - totalSum) === 10 ? 0 : (nextTen - totalSum);
    \end{lstlisting}

    % ==========================================
    \section{Flujo de Registro "Fricción Cero"}
    % ==========================================

    El \texttt{AuthService.register} orquesta una transacción distribuida entre el servicio de identidad de Supabase y las tablas relacionales de la base de datos.

    \subsection{Secuencia de Operaciones}

    \begin{table}[H]
    \centering\small\renewcommand{\arraystretch}{1.5}
    \begin{tabular}{|c|l|p{7cm}|}
    \hline
    \rowcolor{SKTPrimary}
    \textcolor{white}{\textbf{Paso}} & \textcolor{white}{\textbf{Acción}} & \textcolor{white}{\textbf{Descripción}} \\
    \hline
    1 & \textbf{Validación ID} & Ejecución del Módulo 10. Si falla, el flujo termina con Error 400. \\
    \hline
    2 & \textbf{Creación en Auth} & Se crea el usuario en \texttt{supabase.auth} con auto-confirmación de email activada. \\
    \hline
    3 & \textbf{Sincronización Profile} & Inserción en \texttt{public.profiles} con el UUID generado por Auth. \\
    \hline
    4 & \textbf{Sincronización Patient} & Inserción de DNI y metadatos en \texttt{public.patients}. \\
    \hline
    5 & \textbf{Auto-Login} & Se genera la sesión JWT inmediata para que el usuario sea redirigido al Dashboard sin re-autenticarse. \\
    \hline
    \end{tabular}
    \end{table}

    \subsection{Estrategia de Rollback y Robustez}
    Para mantener la integridad referencial, si la inserción en las tablas de base de datos (\texttt{profiles} o \texttt{patients}) falla (por ejemplo, duplicidad de cédula detectada por constraint UNIQUE), el backend solicita inmediatamente el \textbf{borrado del usuario de Auth} mediante el Admin SDK de Supabase, evitando "usuarios huérfanos" en el sistema de seguridad que no existan en el sistema relacional.

    % ==========================================
    \section{Módulo de Sesiones (Login)}
    % ==========================================

    El inicio de sesión delega la autenticación criptográfica a Supabase, pero el backend actúa como guardián de la lógica de negocio.

    \begin{itemize}
        \item \textbf{Carga de Claims}: El JWT resultante contiene el identificador único del usuario (\texttt{sub}) y se orquesta con las políticas RLS para acceso a datos.
        \item \textbf{Manejo de Errores}: Los intentos fallidos son capturados y estandarizados para no revelar información sobre la existencia de correos electrónicos en la base de datos (Seguridad por Ofuscación).
    \end{itemize}

    % ==========================================
    \section{Responsabilidades}
    % ==========================================

    \begin{table}[H]
    \centering\small\renewcommand{\arraystretch}{1.4}
    \begin{tabular}{|p{4.5cm}|p{9cm}|}
    \hline
    \rowcolor{SKTPrimary}
    \textcolor{white}{\textbf{Responsable}} & \textcolor{white}{\textbf{Rol en el Módulo}} \\
    \hline
    Daniel Luisa & Implementación de lógica NestJS, Integración Supabase SDK y validadores. \\
    \hline
    Sebastián Ortiz & Diseño de transaccionalidad, constraints SQL y revisión de seguridad. \\
    \hline
    \end{tabular}
    \end{table}

    % ==========================================
    \section{Firmas de Responsabilidad}
    % ==========================================

    \renewcommand{\arraystretch}{2}
    \noindent
    \begin{tabular}{
    p{0.28\textwidth}
    p{0.30\textwidth}
    >{\centering\arraybackslash}p{0.22\textwidth}
    >{\centering\arraybackslash}p{0.12\textwidth}
    }
    \hline
    \textbf{Rol} &
    \textbf{Nombre / Representante} &
    \textbf{Firma (QR)} &
    \textbf{Fecha} \\
    \hline

    Gestor del Proyecto
    & Angel Ayuquina
    & \includegraphics[width=3.2cm]{QR_Sello_Angel_Ayuquina.png}
    & 07/04/2026 \\[0.6cm]

    Arquitecto BD y Cloud
    & Sebastián Ortiz
    & \includegraphics[width=3.2cm]{QR_Sello_Sebastían_Ortiz.png}
    & 07/04/2026 \\[0.6cm]

    Desarrollador Backend y QA
    & Daniel Luisa
    & \includegraphics[width=3.2cm]{QR_Sello_Daniel_Luisa.png}
    & 07/04/2026 \\[0.6cm]

    Desarrollador Frontend y UX/UI
    & Alex Huachi
    & \includegraphics[width=3.2cm]{QR_Sello_Alex_Huachi.png}
    & 07/04/2026 \\[0.6cm]
    \hline
    \end{tabular}

    % ==========================================
    \section{Control de Cambios}
    % ==========================================

    \renewcommand{\arraystretch}{1.5}

    \noindent
    \begin{tabular}{|p{1.5cm}|p{2.5cm}|p{6.5cm}|p{3.2cm}|}
    \hline
    \rowcolor{SKTLight}
    \textbf{Versión} &
    \textbf{Fecha} &
    \textbf{Descripción del Cambio} &
    \textbf{Responsable} \\
    \hline
    1.0 & 07/04/2026 & Creación del documento de Autenticación y Cédula (Módulo 10). & Daniel Luisa \\
    \hline
    \end{tabular}

    % ==========================================
    \section{Anexos}
    % ==========================================

    \section*{A. Codificación del documento}

    \subsection*{A.1 Definición}
    La codificación es el sistema alfanumérico estructurado que identifica de forma única cada documento dentro del proyecto, permitiendo clasificar, ordenar, recuperar y dar trazabilidad a los documentos.

    \subsection*{A.2 Estructura de codificación}
    La estructura utilizada en el proyecto es la siguiente:
    \[
    \texttt{XX-YYY-ZZZ-NNN}
    \]
    \begin{itemize}
      \item \texttt{XX} — Tipo de proceso (2 letras).
      \item \texttt{YYY} — Nombre del proceso o subproceso (3 letras).
      \item \texttt{ZZZ} — Tipo documental (3 letras).
      \item \texttt{NNN} — Número secuencial (3 dígitos, con relleno a la izquierda).
    \end{itemize}

    \subsection*{A.3 Catálogo de códigos}
    \begin{table}[h!]
    \centering
    \small
    \begin{tabularx}{\linewidth}{@{}lX@{}}
    \toprule
    \textbf{Campo} & \textbf{Opciones (código)} \\ \midrule

    Tipo de proceso (XX) &
    Directivo (DI), Misional (MI), Asesor (AS), Apoyo (AP). \\

    Nombre del proceso (YYY) &
    Desarrollo SW (DES), Gestión de Requisitos (GDR), Implementación del Sistema Clínico (ISC), 
    Gestión de Seguridad de la Información (GSI), Gestión de Agendamiento Médico (GAM), 
    Gestión de Indicaciones Médicas (GIM). \\

    Tipo de documento (ZZZ) &
    Cronograma (CRO), Especificación (ESP), Instructivo (INS), Manual (MAN), Memorando (MEM), 
    Oficio (OFI), Plan (PLA), Procedimiento (PRO), Registro (RGI), Bitácora (BIT), 
    Protocolo (PTC), Guía (GUI). \\

    Número secuencial (NNN) &
    001, 002, … (numeración incremental según la combinación anterior). \\

    \bottomrule
    \end{tabularx}
    \caption*{\footnotesize Listado resumido de códigos importantes.}
    \end{table}

    \subsection*{A.4 Ejemplo de código final}
    Aplicando las selecciones resumidas:
    \begin{itemize}
      \item Tipo de proceso: \textbf{MI} (Misional)
      \item Proceso / Subproceso: \textbf{DES} (Desarrollo SW)
      \item Tipo de documento: \textbf{MAN} (Manual)
      \item Secuencial: \textbf{001}
    \end{itemize}

    \noindent\textbf{Código resultante:} \texttt{MI-DES-MAN-001}

    \subsection*{A.5 Plantilla corta de registro (para control y trazabilidad)}
    Incluye siempre al crear o modificar un documento una entrada como la siguiente en el índice de control:

    \begin{table}[H]
    \centering
    \begin{tabular}{p{3.6cm} p{7cm}}
    \toprule
    \textbf{Campo} & \textbf{Valor / Ejemplo} \\ \midrule
    Código documento & \texttt{MI-DES-MAN-001} \\
    Título & Backend: Autenticación y Cédula \\
    Versión & 1.0 \\
    Fecha emisión / modificación & 2026-04-07 \\
    Autor / Responsable & Daniel Luisa (Backend) \\
    Motivo del cambio & Creación inicial \\
    Aprobación & Angel Ayuquina / 2026-04-07 \\
    Ubicación archivo & \texttt{/docs/MI-DES-MAN-001\_Backend\_Autenticacion\_Cedula.tex} \\
    \bottomrule
    \end{tabular}
    \caption*{\footnotesize{Plantilla mínima para registro de cada documento.}}
    \end{table}

    \subsection*{A.6 Notas sobre uso}
    \begin{itemize}
      \item El número secuencial se incrementa para cada nuevo documento dentro de la misma combinación \texttt{XX-YYY-ZZZ}.  
      \item Mantener el registro de cambios vinculado al repositorio o al sistema documental para garantizar trazabilidad.  
    \end{itemize}

    \section*{B. Paleta de colores y archivos fuente}

    \subsection*{B.1 Introducción}
    Este anexo contiene la paleta de color muestreada a partir del archivo raster recibido, la especificación técnica básica (HEX / RGB / CMYK), y resultados de contraste WCAG frente a blanco y negro.

    \subsection*{B.2 Imagen fuente}
    \begin{figure}[H]
      \centering
      \includegraphics[width=0.95\linewidth]{SKTLogo.png}
      \caption*{Figura 1: Logo del equipo SKT Software Solution}
    \end{figure}

    \subsection*{B.3 Tabla técnica — paleta}
    \small
    \begin{table}[H]
    \centering
    \begin{tabular}{ l  c  c  c  c }
    \toprule
    \textbf{Nombre} & \textbf{Muestra} & \textbf{HEX} & \textbf{RGB} & \textbf{CMYK (\%)} \\
    \midrule
    Primario (sugerido) & \cellcolor{SKTPrimary} & \texttt{\#3F4A4F} & (63,74,79) & (20,6,0,69) \\
    Secundario (oscuro) & \cellcolor{SKTDark} & \texttt{\#242B2E} & (36,43,46) & (22,7,0,82) \\
    Matiz medio & \cellcolor{SKTMid} & \texttt{\#6E7375} & (110,115,117) & (6,2,0,54) \\
    Acento & \cellcolor{SKTAccent} & \texttt{\#979C9D} & (151,156,157) & (4,1,0,38) \\
    Claro (detalle) & \cellcolor{SKTLight} & \texttt{\#CDCFD2} & (205,207,210) & (2,1,0,18) \\
    \bottomrule
    \end{tabular}
    \end{table}
    \normalsize

    \subsection*{B.4 Interpretación de contraste WCAG}
    \begin{itemize}
      \item Ratios mostrados siguen la fórmula WCAG (relativa a luminancias sRGB linealizadas).
      \item \textbf{Regla práctica:} para texto normal, se recomienda contraste \(\geq 4.5:1\) (AA). Para texto grande, \(\geq 3:1\).
      \item De la paleta muestreada: \textbf{texto blanco sobre el color Primario (\#3F4A4F)} cumple AA (9.12), mientras que \textbf{texto negro sobre ese mismo fondo} no cumple (2.30). Eso implica que en fondos oscuros del isotipo es preferible usar texto/elementos en blanco o variantes claras.
      \item Colores muy claros como \#CDCFD2 no son adecuados para texto claro encima; verificar siempre combinaciones en contexto.
    \end{itemize}

    \end{document}